% ---------------------------------------------------------------------------
%  03-algorithm.tex -- turning the theorem into a test
% ---------------------------------------------------------------------------
\section{The algorithm}
\label{sec:algorithm}

\subsection{Partitioning a node}

\code{components()} treats each goal of $\Gam$, each context formula of $\Ctx$
and each live pursuit as one \emph{item}, and runs union--find over the items,
joining two whenever their symbol sets meet. That is \ref{S1}.

Atom sets are \textbf{bitmasks}, one bit per atom of the input, memoised per
formula id. The test compares every pair of items at every node, so as sets of
strings it is three nested loops with a string compare at the bottom, and it
measured about a third of the run on formulas that never split. As words it is
one \code{\&}. An input with more than $64$ atoms simply does not decompose.

\subsection{The side conditions, and where each bites}

\paragraph{Clock ownership (\ref{S2}).} Ownership is read off the
\emph{formulas}, not off the live pursuits, via \code{clocks\_of()}: a modality
still sitting in $\Gam$ already owns its formula-clock. Deriving ownership from
pursuits alone leaves such a clock owned by nobody, and the component
construction then frees it out of \emph{every} component's zone. That is a
weakening, and it is the exact bug that made
$\mathrm{AG}(\mathrm{EF}_{[2,4]}p)\wedge\mathrm{AG}_{[4,7]}\neg p$ come back
satisfiable: freeing the onset clock erases the constraint that distinguishes
one onset from another. A clock owned by more than one group, or constrained
but owned by none, aborts the split.

\paragraph{Shared symbols are not only propositions.} A data variable couples
its constraints exactly as an atom does: $\{x>3\}$ and $\{x<2\}$ mention no
proposition at all, and grouping by propositions alone puts them in different
components, each satisfiable on its own. \code{IntVar} names therefore join the
bitmask alongside \code{Atom} names. The self-test caught this on twelve checks
of the form $\{x\le 5\}\wedge\{x\ge 10\}$; the sweeps did not, having no data
variables in them.

\paragraph{Degenerate components.} A component mentioning no atom at all --- a
bare $\top$ or $\bot$ --- is folded into the first real component instead of
being split off. Peeling those off is all cost and no benefit, and empirically
it was also where the first version went wrong: every split it performed on
$\mathrm{AG}(\mathrm{EF}_{[2,4]}p)\wedge\mathrm{AG}_{[4,7]}\neg p$ was
degenerate, and the verdict was still wrong, which is what localised the fault
to clock ownership rather than to the partition.

\subsection{Deciding by components}

\code{decide\_by\_components()} keys each component on its node key and looks it
up in a cache shared with every sub-search. On a miss it runs a fresh engine
from that node --- constructed from the original root formula, so the clock
allocation, the maximum constants and the split points are identical --- and
stores the verdict. One false component is enough to reject.

The cache is what makes this component \emph{caching} rather than component
\emph{splitting}: the same component recurs across many parent nodes, which is
precisely why the undecomposed search enumerated their product.

\subsection{Hooking it into the search}

\code{try\_split} runs in \code{decide\_scc} before a node is first visited. A
node that splits is marked settled with the conjunction of its components'
verdicts and is never expanded: it becomes a leaf of the enclosing search,
exactly like a node whose own SCC has already closed, which the search already
knows how to consume. Tarjan never sees it, so no cycle can run through it ---
which is sound because \cref{cor:search} decides it outright.

Recursion terminates because every component is strictly smaller than the node
it came from; a depth cap guards against pathology.

\subsection{Not re-deciding shapes that never split}

Whether a node splits depends only on \emph{which formulas it carries}, never on
its zone. So a shape that did not split never will, and \code{nosplit\_} records
its signature. Shapes repeat heavily --- a nested tower has roughly half as many
shapes as nodes --- and without the memo the union--find runs again at each of
them, which is pure loss on any formula that never decomposes. This is what
brings the overhead on non-splitting inputs down to a few per cent
(\cref{sec:eval}).

\subsection{Two closures that were buying nothing}

Two costs found by re-profiling after the decomposition work, both unrelated to
it and both removed on a proof rather than a measurement.

\emph{The dead-clock closure.} Eliminating the clocks a node no longer reads
used to re-close the DBM, which was $30\%$ of the run. It is unnecessary. The
update only relaxes, and the result is already canonical: checking
$d[i][j]\le d[i][l]+d[l][j]$ on the new matrix given it held on the old one, a
path through a freed clock is $d[i][0]+\infty=\infty$; a freed row is $\infty$
throughout; a freed column becomes $d[i][0]$, and old closure already gave
$d[i][0]\le d[i][l]+d[l][0]$. The reference row needs
$d[0][l]+d[l][0]\ge 0$, which holds because a closed non-empty zone has no
negative cycle.

\emph{The extrapolation closure.} $\mathrm{Extra}_M$ closed unconditionally; a
zone none of whose bounds exceeded its clock's maximum is unchanged and arrived
closed, so a dirty flag skips the cubic step.

\subsection{The static form}

The same theorem applied once, at the root, to the top-level conjunction is a
preprocessing step: split the conjuncts into atom-disjoint groups, solve each,
conjoin. It costs one test and it is strictly stronger on the shapes it catches,
because no per-node bookkeeping is paid at all. It catches nothing whose
independence only emerges during the search. Both forms are implemented; see
\cref{sec:eval} for the comparison.
